\section{Numerical precision is part of a financial method}
\label{sec:numerical-methods}

The affine family is useful for Northstar because the scenario fixes prices,
rates and operating coefficients while the debt share remains unknown. Other
contracts require roots or discount factors. Their numerical treatment must
state what is exact, what is enclosed and when a posted amount is determined.
An accounting name on a library call does not make its result exact.

Consider positive payments $C_i$ at positive times $t_i$ and a positive price
$P$. The effective annual rate solves
\[
 f(r)=\sum_i\frac{C_i}{(1+r)^{t_i}}-P=0,\qquad r>-1.
\]
Each term discounts a payment to the valuation date. In this positive-payment
case,
\[
 f'(r)=-\sum_i\frac{t_i C_i}{(1+r)^{t_i+1}}<0.
\]
Consequently a certified sign-changing bracket contains exactly one root.
The method uses stated bracket bounds and an expansion cap, evaluates signs
with outward error bounds, and bisects until the requested rate or value
enclosure is reached. An uncertified sign, unavailable bracket or exhausted
accuracy budget yields an unresolved numerical result. Mixed-sign cash flows
can have several internal rates of return; the preceding monotonicity argument
does not apply and cannot justify returning the first Newton iterate that
converges.

Net proceeds 100 and a payment 121 in two years give exactly 10\%. An immediate
fee of 2 changes proceeds to 98 and changes that root. Annual-effective 21\%
for half a year has the exact accumulation factor $\sqrt{1.21}=1.1$, so interest
on 100 is 10. Simple proration would give 10.5 and computes a different contract.
At annual-effective 10\%, the half-year factor $\sqrt{1.1}$ is irrational. For
rational time $p/q$ and positive $1+r$, exact rational power comparisons can
enclose the positive $q$th root of $(1+r)^p$. A negative $p$ first requires the
reciprocal. General exponential discounting needs a separately specified
error-controlled method.

A posting quantum of 0.01 does not eliminate this uncertainty. The whole
certified monetary enclosure must lie in one rounding cell under the admitted
tie rule. An enclosure from 1.0049 to 1.0051 crosses a half-cent boundary and
does not establish one posted amount. Refine within a declared budget or retain
rounding unresolved. This rule separates evidence for a numerical answer from
the accounting authority to use that answer.

\section{Updating a law without counting evidence twice}
\label{sec:inference}

Two financial histories can be compatible with the same reports without being
equally probable. Once an explicitly justified law is supplied, genuinely new
evidence may change it. For finite histories $h$, prior mass $p(h)$ and an
admitted observation likelihood $L(D\mid h)$, Bayes' rule gives
\[
 p(h\mid D)=\frac{p(h)L(D\mid h)}{\sum_g p(g)L(D\mid g)}.
\]
The denominator is the probability of that evidence under the declared model.
If it is zero, the condition is impossible in that model; dropping incompatible
histories and renormalising an empty collection cannot repair it.

For two histories with assumed prior masses $1/3$ and $2/3$ and likelihoods
$3/4$ and $1$, the unnormalised masses are $1/4$ and $2/3$. Their sum is $11/12$,
so the posterior masses are $3/11$ and $8/11$. If the second history's cash
calculation fails, the first still has mass $3/11$ and unresolved cash retains
mass $8/11$. A report giving the successful history probability one would
answer a different, success-conditioned question.

An observation has a semantic information identity as well as a file locator.
Reimporting the same observation through another file does not create a second
measurement. Repeating an identical admitted condition is idempotent: it keeps
the law and information lineage unchanged. Submitting an already consumed
observation as a fresh likelihood block is rejected as duplicate information.
Equal selected member sets alone do not prove that two conditions have the same
meaning. An admitted conditional kernel may divide a region's existing mass
when it is refined; the number of new representatives never determines that
division.

Continuous inference adds a measure choice. One must specify coordinates $z$,
domain, base measure, prior, observation mechanism and any transformation
Jacobian. \citet[Section 2.2, equations (2.5)--(2.8)]{stuart2010} supplies the
inverse-problem formulation, but does not establish identification or easy
sampling for a mixed accounting model. In particular, an exact accounting
equality can confine a model to a lower-dimensional set. Multiplying an ordinary
Lebesgue density by its zero-volume indicator is not a posterior construction.
One must parameterise the constrained model with its declared measure or supply
an actual measurement-error model. Adding convenient small noise changes the
target.

\section{A continuous example whose posterior can be checked}

A linear Gaussian model supplies a tractable comparison for later numerical
inference. Assume the unknown economic quantity has a Gaussian prior
$z\sim\mathcal{N}(m_0,V_0)$ with positive definite covariance $V_0$.
Let the observation be $y=Hz+\epsilon$, where
$\epsilon$ is independent Gaussian noise with positive definite covariance
$R$. The measurement-error assumption is substantive: it describes this example,
not rounding in every financial statement.

Multiplying prior and likelihood adds their quadratic forms. Terms that depend
on $z$ are
\begin{align*}
 &(z-m_0)^\top V_0^{-1}(z-m_0)
 +(y-Hz)^\top R^{-1}(y-Hz)\\
 &\quad=z^\top(V_0^{-1}+H^\top R^{-1}H)z
 -2z^\top(V_0^{-1}m_0+H^\top R^{-1}y)+\text{constant}.
\end{align*}
Completing the square therefore gives posterior covariance $V$ and mean $m$
through
\[
 V^{-1}=V_0^{-1}+H^\top R^{-1}H,\qquad
 V^{-1}m=V_0^{-1}m_0+H^\top R^{-1}y.
\]
The prior and observation contribute precision and precision-weighted
information. An implementation should solve factored linear systems rather
than explicitly form these inverse matrices, and reject an invalid covariance
or unsupported rank condition.

For $m_0=0$, $V_0=4$, $H=1$, $R=1$ and $y=2$, posterior precision is
$1/4+1=5/4$, variance is $4/5$ and mean is $(4/5)2=8/5$. A conditional cash
quantity $100+10z$ has mean 116 and variance 80. A second \emph{independent}
measurement of 2 would give precision $9/4$, variance $4/9$ and mean $16/9$.
Reusing the first measurement would produce those same numbers but would be
wrong for the stated information. The analytic example can distinguish numerical
error from information double counting.

For a bounded non-Gaussian scalar model, quadrature requires a stated domain
or tail truncation and sufficient integrand error bounds to enclose the
normalizer and requested moments. Without a positive lower normalizer bound,
the ratio remains unresolved. Mixed models can enumerate finite choices and
integrate each conditional continuous model, but combining branches requires
their evidence factors, not draw counts. Optional MCMC needs target-specific
parameterisation, tuning, independent chains, tractable correctness comparisons,
convergence and Monte Carlo error checks. This proposal supplies no generally
validated sampler or universal tolerance.

Nor can a more elaborate inference engine identify a quantity that the
observation model never sees. If reports constrain only total debt 100, the
likelihood may be constant in its floating share $x$. Its conditional prior
then remains unchanged. A uniform prior is a chosen assumption, not a discovery
from those reports. The inspected working version of
\citet[Sections 2.1--2.3, equations (2.2)--(2.4)]{giacomini2021} motivates
retaining posterior classes where a conditional prior cannot be revised.
For Northstar there is a simple direct deduction: under all probability
measures on $[0,1]$, the expectation of $5+2x$ lies in $[5,7]$, with endpoints
attained by point masses. This deduction imports no asymptotic theorem from
that paper's structural-vector-autoregression setting.

\section{A stress is not an observed condition or an intervention}

The interest-rate stress in Northstar fixes a future driver path. It supplies
no new evidence about the historical debt share. An observational forecast
instead conditions a joint law on information. A causal intervention needs
structural equations and an exogenous disturbance law, because it replaces
the mechanism determining a chosen variable while leaving other mechanisms
in place \citep[Section 3.2, equations (5)--(9) and (15)--(17)]{pearl2009}.

An elementary example makes the difference numerical. Let $U$ equal zero or
one with probability $1/2$ each, and assume
\[
       X=U,\qquad Y=2X+U.
\]
Observing $X=1$ establishes $U=1$ and hence $Y=3$. Intervening to set $X=1$
replaces the first equation. It leaves the original law of $U$, so $Y$ is 2 or
3 with equal mass and mean 2.5. Selection of a high-investment company and an
imposed investment policy can therefore answer different questions even under
one fully specified model. Financial filings alone do not establish these
causal equations or identify their effects.

An intervention request must name its target, replacement equation, structural
parents and disturbance law and satisfy acyclicity or a separately supported
equilibrium rule. Correlated disturbances remain joint. The initial company
profile supports imposed exogenous stress; the continuous and causal methods
described here require separately admitted extension schemas and signatures.
They do not silently expand the initial finite-law interface.

At every decision date, behavior and measurement receive the admitted historical
view and permitted simulated past. The scheduler's knowledge of a complete
future stress path is not permission for a manager to foresee it. A perfect-
foresight policy needs an explicit additional assumption. Common random numbers
can improve a paired numerical comparison while establishing no causal effect
in actual company data.
